\scanpage{108}
\mspage{53}
canonique
\[
D_{\deleted{\unsure{p}}d}\longrightarrow X
\]
\margin{On introduit aussi $D_0=X$, (152) et $D_0\longrightarrow X$ l'identité}
qui se factorise d'ailleurs par $D$, et qui est un morphisme d'immersion local, de codim $p$. Pour $D_{\deleted{\unsure{p}}d}$, on a un revêt\textsuperscript{t} étale \unsure{à|:} $p$ feuillets, noté $\overset{\deleted{\unsure{q}}}{\widetilde D_{\deleted{\unsure{p}}d}}$, \deleted{formé des couples}\inserted{dont la fibre en} un pt $(x,V)$ est \unsure{formée|formé} de l'ens des $\deleted{\unsure{p}}d$ branches $D_{xi}$ de $D$ en $x$ qui contiennent $V$. Donc les \deleted{\illegible{1}} pts de $\widetilde D_{\deleted{\unsure{p}}d}$ correspondent aux triples $(x,V,W)$, où $x$ est un pt de $D$, et $V\subset W$ des germes de ss-multiplicités de $X$ en $x$ attachées à $D$, de codim. respectives $\deleted{\unsure{p}}d$ et $1$. On trouve donc un morphisme canonique
\diagram{\begin{tikzcd}[ampersand replacement=\&, row sep=small]
\widetilde D_{\deleted{\unsure{p}}d}\arrow[r] \& D_1\\
(x,V,W)\arrow[r,mapsto] \& W
\end{tikzcd}}

Plus généralement, si $J$ est un ens. \unsure{\deleted{fini}|fini} ordonné comme ci-dessus, \inserted{muni d'une fonction codim,} on \deleted{\unsure{en}} désigne par $D_J$ la \deleted{\unsure{var.|ens}} multiplicité \deleted{formée}\inserted{dont} les points sont les couples d'un $x\in\deleted{D}\inserted{X}$ \inserted{de codim $d=d(J)$,} et d'un plong\textsuperscript{t} d'ens. ordonnés, compatible \unsure{\deleted{avec}|aux} fonctions codim
\margin{on suppose la fonction codim strictement décr. cf (151)}
\[
(152)\qquad J\mathrel{\overset{\varphi}{\hookrightarrow}}\Delta_x\simeq\mathcal P(I(x))^\circ
\]
\margin{107}
